\scanpage{6}
L'inclusion des parties paraboliques (ou encore, des grpes paraboliques correspondants, \unsure{quand|qnd} le syst. de racines $\Sigma$ provient d'un groupe alg) est alors transformée \deleted{par} \inserted{en l'opposée de} l'inclusion des facettes dans $\check V$. C'est cette relation (\uline{opposée} de l'inclusion des paraboliques) qu'on prend sur cet \deleted{\unsure{e}} ens de parties, qui devient dès lors une géométrie de \uline{drapeaux stricts}, correspondant à un \unsure{squelette|complexe} simplicial, dont l'ens des sommets \inserted{(ou simplexes minimaux)} correspond aux \textquotedblleft paraboliques \unsure{maximaux|minimaux}\textquotedblright, et l'ens des simplexes maximaux aux \textquotedblleft Borels\textquotedblright\ ou \textquotedblleft syst simples de racines\textquotedblright\ ou repères $r\in R$. En fait, il s'agit manifestement d'une triangulation de la sphère \inserted{$S^{l-1}$}. L'ens des repères s'identifiera à $R\times\deleted{\unsure{\mathrm{Ord}}}\inserted{\operatorname{Rep}(I)}$, où $\operatorname{Rep}(I)=\operatorname{Bij}([1,l],I)$ est l'ens des ordres totaux sur l'ens $I$ (de \unsure{cart.|cnt.|crt.|card.} $l$).
% archive sheet number: 5

